\documentclass[12pt,a4paper]{article}

\usepackage[utf8]{inputenc}
\usepackage[T1]{fontenc}
\usepackage[english]{babel}
\usepackage{amsmath}
\usepackage{amssymb}
\usepackage{geometry}
\usepackage{booktabs}
\usepackage{array}
\usepackage{xcolor}
\usepackage{mdframed}
\usepackage{enumitem}
\usepackage{chemformula}

\geometry{margin=2.5cm}

\definecolor{azul}{RGB}{30, 100, 180}
\definecolor{cinza}{RGB}{245, 245, 245}
\definecolor{verde}{RGB}{20, 130, 80}

\mdfdefinestyle{enunciado}{
  backgroundcolor=cinza,
  linecolor=azul,
  linewidth=1pt,
  topline=false,
  rightline=false,
  bottomline=false,
  innerleftmargin=12pt,
  innerrightmargin=12pt,
  innertopmargin=10pt,
  innerbottommargin=10pt
}

\mdfdefinestyle{dados}{
  backgroundcolor=cinza,
  linecolor=gray!40,
  linewidth=0.5pt,
  innerleftmargin=12pt,
  innerrightmargin=12pt,
  innertopmargin=8pt,
  innerbottommargin=8pt,
  roundcorner=4pt
}

\mdfdefinestyle{resolucao}{
  backgroundcolor=white,
  linecolor=verde,
  linewidth=1pt,
  topline=false,
  rightline=false,
  bottomline=false,
  innerleftmargin=12pt,
  innerrightmargin=12pt,
  innertopmargin=10pt,
  innerbottommargin=10pt
}

\begin{document}

\begin{center}
  {\small\textcolor{azul}{\textbf{Física e Química A --- 11.º ano}}}
\end{center}

\bigskip

%---------------------------------------------------------------
% ENUNCIADO
%---------------------------------------------------------------
\begin{mdframed}[style=enunciado]

Numa experiência laboratorial, um aluno pretende determinar a solubilidade do
carbonato de cálcio (\ch{CaCO3}) numa solução aquosa a \(25\,^\circ\text{C}\),
usando um feixe de luz \textit{laser}.

\medskip

O aluno dirige o laser, no ar, para a superfície da solução com um ângulo de
incidência de \(\theta_1 = 45{,}0^\circ\) e mede um ângulo de refração de
\(\theta_2 = 32{,}1^\circ\). Sabe que o índice de refração de uma solução
aquosa aumenta proporcionalmente com a concentração de soluto, de acordo com:

\[
  n_{\text{solução}} = n_{\text{água}} + k \cdot c
\]

sendo \(c\) a concentração em \(\text{mol\,L}^{-1}\) e
\(k = 1{,}45 \times 10^{-2}\,\text{L\,mol}^{-1}\).

\medskip

Sabendo que o \ch{CaCO3} se dissolve segundo o equilíbrio:

\[
  \ch{CaCO3}_{(s)} \rightleftharpoons \ch{Ca^{2+}}_{(aq)} + \ch{CO3^{2-}}_{(aq)}
\]

\end{mdframed}

\medskip

%---------------------------------------------------------------
% DADOS
%---------------------------------------------------------------
\begin{mdframed}[style=dados]
\textbf{Dados:}
\begin{itemize}[leftmargin=1.5em, itemsep=2pt, topsep=4pt]
  \item $K_s\,(\ch{CaCO3},\; 25\,^\circ\text{C}) = 4{,}96 \times 10^{-9}$
  \item $n_{\text{ar}} = 1{,}000$ \quad $n_{\text{água pura}} = 1{,}333$
  \item $M(\ch{CaCO3}) = 100{,}1\,\text{g\,mol}^{-1}$
\end{itemize}
\end{mdframed}

\medskip

%---------------------------------------------------------------
% ALÍNEAS
%---------------------------------------------------------------
\begin{enumerate}[label=\textbf{\alph*)}, leftmargin=1.8em, itemsep=6pt]
  \item Determine o índice de refração da solução.
  \item Calcule a concentração total de iões em solução.
  \item Determine a solubilidade molar do \ch{CaCO3} e compare o valor obtido
        com o calculado a partir de $K_s$. Os resultados são consistentes?
\end{enumerate}

\bigskip\bigskip

%---------------------------------------------------------------
% RESOLUÇÃO
%---------------------------------------------------------------
\noindent\textcolor{verde}{\rule{\linewidth}{0.5pt}}

\medskip
\noindent{\textcolor{verde}{\textbf{Resolução}}}
\medskip

\noindent\textcolor{verde}{\rule{\linewidth}{0.5pt}}

\bigskip

%--- alínea a)
\begin{mdframed}[style=resolucao]
\textbf{a) Índice de refração da solução --- Lei de Snell-Descartes}

\medskip

Aplicando a lei de Snell-Descartes na interface ar/solução:

\[
  n_{\text{ar}}\,\sin\theta_1 = n_{\text{sol}}\,\sin\theta_2
\]

\[
  1{,}000 \times \sin(45{,}0^\circ) = n_{\text{sol}} \times \sin(32{,}1^\circ)
\]

\[
  n_{\text{sol}} = \frac{0{,}7071}{0{,}5314} \approx \boxed{1{,}331}
\]
\end{mdframed}

\bigskip

%--- alínea b)
\begin{mdframed}[style=resolucao]
\textbf{b) Concentração total de iões}

\medskip

Isolando \(c\) na expressão do índice de refração:

\[
  c = \frac{n_{\text{sol}} - n_{\text{água}}}{k}
\]

Usando o valor mais preciso \(n_{\text{sol}} = 1{,}3331\):

\[
  c = \frac{1{,}3331 - 1{,}3330}{1{,}45 \times 10^{-2}}
    \approx \boxed{6{,}9 \times 10^{-5}\,\text{mol\,L}^{-1}}
\]

\medskip
Como \([\ch{Ca^{2+}}] = [\ch{CO3^{2-}}] = s\), a concentração total de iões
é \(2s \approx 6{,}9 \times 10^{-5}\,\text{mol\,L}^{-1}\).
\end{mdframed}

\bigskip

%--- alínea c)
\begin{mdframed}[style=resolucao]
\textbf{c) Solubilidade molar e comparação com $K_s$}

\medskip

\textit{Valor experimental} (da alínea b):

\[
  s_{\text{exp}} = \frac{c_{\text{total}}}{2} \approx 3{,}45 \times 10^{-5}\,\text{mol\,L}^{-1}
\]

\textit{Valor teórico} a partir de $K_s$:

\[
  K_s = [\ch{Ca^{2+}}][\ch{CO3^{2-}}] = s^2
  \implies s = \sqrt{K_s} = \sqrt{4{,}96 \times 10^{-9}}
  \approx \boxed{7{,}04 \times 10^{-5}\,\text{mol\,L}^{-1}}
\]

\medskip

Os dois valores são da mesma ordem de grandeza ($10^{-5}\,\text{mol\,L}^{-1}$),
confirmando a consistência do método. A pequena diferença deve-se ao
arredondamento nos ângulos medidos --- uma variação de décimas de grau numa
solução tão diluída traduz-se numa grande incerteza no índice de refração.
\end{mdframed}

\end{document}
